\section{Примеры}

В этом разделе собраны примеры, на которых видны все понятия главы. Сначала~--- необходимые сведения о $p$-адических
числах.

\subsection*{$p$-адические числа}

\begin{example}[7.33: кольцо $\Zp$ и поле $\Qp$]
Пусть $p$~--- простое число. Кольцо \emph{$p$-адических целых чисел}~--- это $\Zp=\varlim_k\Z/p^k\Z$ (с редукциями
$\Z/p^{k+1}\to\Z/p^k$), т.\,е.\ согласованные последовательности вычетов по модулям $p,p^2,p^3,\dots$
\begin{itemize}[leftmargin=*]
\item Каждый элемент однозначно записывается рядом $\alpha=a_0+a_1p+a_2p^2+\cdots$ с <<цифрами>> $a_i\in\{0,\dots,p-1\}$,
понимаемым так: $\alpha\bmod p^k=a_0+\dots+a_{k-1}p^{k-1}$. Сложение и умножение~--- столбиком с переносом в следующий
разряд. Например, $-1=(p-1)+(p-1)p+(p-1)p^2+\cdots$ (прибавьте $1$: все разряды обнулятся с переносом), а
$\frac1{1-p}=1+p+p^2+\cdots$. Вообще, каждое целое $m$ с $p\nmid m$ обратимо в $\Zp$ (оно обратимо по модулю каждого $p^k$).
\item \emph{Единицы:} $\Zp^\times=\{\alpha\mid a_0\ne0\}$. Каждый ненулевой $\alpha$ однозначно записывается как $p^vu$,
$v\ge0$, $u\in\Zp^\times$; число $v=v_p(\alpha)$~--- \emph{$p$-адическое нормирование}. Отсюда $\Zp$~--- область
целостности с единственным максимальным идеалом $p\Zp$ и $\Zp/p\Zp=\F_p$.
\item \emph{Метрика:} $|\alpha|_p=p^{-v_p(\alpha)}$, $|0|_p=0$; $|\alpha\beta|_p=|\alpha|_p|\beta|_p$ и
$|\alpha+\beta|_p\le\max(|\alpha|_p,|\beta|_p)$. Шары~--- смежные классы $a+p^k\Zp$, так что метрическая топология
совпадает с топологией предела; $\Zp$ компактно, вполне несвязно, и $\Z$ плотно в $\Zp$.
\item \emph{Поле $p$-адических чисел} $\Qp=\Zp[1/p]$~--- поле частных $\Zp$, его элементы~--- ряды $\sum_{i\ge-m}a_ip^i$.
Это пополнение $\Q$ относительно $|\cdot|_p$ (как $\R$~--- пополнение $\Q$ относительно обычного модуля).
\item \emph{Связь с $\Zh$.} По китайской теореме $\Z/n\cong\prod_{\ell^k\| n}\Z/\ell^k$ согласованно по $n$, поэтому
$\Zh\cong\prod_{\ell}\Z_\ell$ (произведение по всем простым $\ell$), как кольца и как топологические группы.
\end{itemize}
\end{example}

\begin{center}
\begin{tikzpicture}[level distance=1.05cm,
  level 1/.style={sibling distance=6.4cm},
  level 2/.style={sibling distance=3.2cm},
  level 3/.style={sibling distance=1.6cm},
  every node/.style={font=\footnotesize}, edge from parent/.style={draw,gray}]
\node {$\Z_2$}
  child {node {$0$}
    child {node {$0$} child {node {$0$}} child {node {$4$}}}
    child {node {$2$} child {node {$2$}} child {node {$6$}}}}
  child {node[cnew] {$1$}
    child {node {$1$} child {node {$1$}} child {node {$5$}}}
    child {node[cnew] {$3$} child {node {$3$}} child {node[cnew] {$7$} edge from parent[draw=cnew,thick]}
      edge from parent[draw=cnew,thick]}
    edge from parent[draw=cnew,thick]};
\node at (-7.1,-1.05) {$\Z/2$};
\node at (-7.1,-2.1) {$\Z/4$};
\node at (-7.1,-3.15) {$\Z/8$};
\node at (-7.1,-3.9) {$\vdots$};
\end{tikzpicture}
\end{center}
\figcap{$\Z_2=\varlim\Z/2^k$ как дерево: вершины уровня $k$~--- вычеты по модулю $2^k$, рёбра~--- редукция. Точка $\Z_2$~---
бесконечный путь вниз (красный путь $1,3,7,\dots$~--- это $-1=1+2+4+\cdots$); базисная окрестность $a+2^k\Z_2$~--- все
пути через вершину $a$ уровня $k$. Два разных пути рано или поздно расходятся~--- отсюда хаусдорфовость и вполне
несвязность; компактность~--- из того, что на каждом уровне вершин конечное число.}

\begin{lemma}[7.33a: Гензель]
Пусть $f\in\Zp[X]$, $a_0\in\Z$, $f(a_0)\equiv0\pmod p$ и $f'(a_0)\not\equiv0\pmod p$. Тогда существует единственный
$a\in\Zp$ с $f(a)=0$ и $a\equiv a_0\pmod p$.
\end{lemma}

\begin{proof}
Пусть уже найден $a_k$ с $f(a_k)\equiv0\pmod{p^k}$, $a_k\equiv a_0\pmod p$. Ищем $a_{k+1}=a_k+tp^k$: по формуле Тейлора
$f(a_{k+1})\equiv f(a_k)+tp^kf'(a_k)\pmod{p^{k+1}}$, и так как $f'(a_k)\equiv f'(a_0)\not\equiv0\pmod p$, вычет
$t\bmod p$ определяется однозначно. Согласованная последовательность $(a_k)$~--- искомый элемент $\Zp$; единственность
следует из однозначности $t$ на каждом шаге.
\end{proof}

Например, $X^2+1$ имеет корень в $\Z_5$ ($2^2\equiv-1\pmod5$), но не в $\Z_3$ ($-1$~--- не квадрат по модулю 3). Многочлен
$X^{p-1}-1$ раскладывается над $\F_p$ на различные линейные множители, поэтому $\Zp$ содержит группу $\mu_{p-1}$ всех
корней степени $p-1$ из единицы, и редукция $\mu_{p-1}\to\F_p^\times$~--- изоморфизм.

\begin{proposition}[7.33b: строение $\Zp^\times$]
\begin{enumerate}[label=\textup{(\alph*)}]
\item При нечётном $p$: $\Zp^\times=\mu_{p-1}\times(1+p\Zp)$, и $a\mapsto(1+p)^a$~--- изоморфизм топологических групп
$\Zp\to1+p\Zp$. Итого $\Zp^\times\cong\Z/(p-1)\times\Zp$.
\item $\Z_2^\times=\{\pm1\}\times(1+4\Z_2)$, и $a\mapsto5^a$~--- изоморфизм топологических групп $\Z_2\to1+4\Z_2$. Итого
$\Z_2^\times\cong\Z/2\times\Z_2$.
\end{enumerate}
\end{proposition}

\begin{proof}
(a) Редукция $\Zp^\times\to\F_p^\times$ сюръективна с ядром $1+p\Zp$, а её ограничение на $\mu_{p-1}$~--- изоморфизм;
отсюда разложение в прямое произведение. Далее, при нечётном $p$
\[
(1+p)^{p^k}\equiv1+p^{k+1}\pmod{p^{k+2}}\qquad(k\ge0).
\]
Индукция: если $x=1+p^{k+1}+cp^{k+2}$, то $x^p=1+p(p^{k+1}+cp^{k+2})+\binom p2(p^{k+1}+\dots)^2+\cdots\equiv1+p^{k+2}
\pmod{p^{k+3}}$, так как $\binom p2$ делится на $p$ (здесь нужна нечётность $p$), а остальные члены делятся на
$p^{3(k+1)}$. Значит, порядок $1+p$ в $(\Z/p^{k+1})^\times$ равен ровно $p^k$. Подгруппа $U_{k+1}=\ker\bigl((\Z/p^{k+1})^\times\to\F_p^\times\bigr)$
(вычеты $\equiv1\pmod p$) имеет $p^k$ элементов и, следовательно, циклична с образующей $1+p$. Изоморфизмы
$\Z/p^k\to U_{k+1}$, $a\mapsto(1+p)^a$, согласованы по $k$, а $\varlim_kU_{k+1}=1+p\Zp$, и в пределе дают непрерывный изоморфизм $\Zp\to1+p\Zp$ (этим и определяется
$(1+p)^a$ для $a\in\Zp$); он гомеоморфизм по компактности.

(b) Редукция $\Z_2^\times\to(\Z/4)^\times=\{\pm1\}$ имеет ядро $1+4\Z_2$ и расщепляется ($\pm1\in\Z_2^\times$). Далее,
$5^{2^k}\equiv1+2^{k+2}\pmod{2^{k+3}}$: если $x=1+2^{k+2}+c2^{k+3}$, то $x^2=1+2(2^{k+2}+c2^{k+3})+(2^{k+2}+\dots)^2\equiv
1+2^{k+3}\pmod{2^{k+4}}$, так как $2(k+2)\ge k+4$. Значит, порядок $5$ по модулю $2^{k+2}$ равен $2^k$ (например, по
модулю $64$ он равен $16$), и, как в (a), $a\mapsto5^a$~--- изоморфизм $\Z_2\to1+4\Z_2$. (Группа $1+2\Z_2=\Z_2^\times$
не проциклическая: она содержит элемент $-1$ порядка 2, и, например, $3^2\equiv1\pmod8$.)
\end{proof}

\subsection*{Замкнутые подгруппы и подполя $\Fpb$}

\begin{proposition}[7.34]
\begin{enumerate}[label=\textup{(\alph*)}]
\item Замкнутые подгруппы $\Zp$~--- это $0$ и $p^n\Zp$, $n\ge0$.
\item Замкнутые подгруппы $\Zh=\prod_\ell\Z_\ell$~--- это произведения $\prod_\ell H_\ell$, где каждое
$H_\ell\in\{0\}\cup\{\ell^n\Z_\ell\mid n\ge0\}$.
\item Каждая подгруппа конечного индекса в $\Zh$ открыта и имеет вид $n\Zh$.
\end{enumerate}
\end{proposition}

\begin{proof}
(a) Пусть $H\ne0$ замкнута. Для $h\in H$ все $mh$, $m\in\Z$, лежат в $H$, а $\Z$ плотно в $\Zp$, умножение непрерывно и
$H$ замкнута; поэтому $\Zp h\subset H$. Выберем $h\in H$ с наименьшим $v_p(h)=n$: $h=p^nu$, $u\in\Zp^\times$, и
$p^n\Zp=\Zp h\subset H\subset p^n\Zp$.

(b) По тому же рассуждению замкнутая $H\subset\Zh$~--- идеал. Пусть $e_\ell\in\Zh$~--- идемпотент с компонентой $1$ в
$\Z_\ell$ и $0$ в остальных. Тогда $e_\ell H\subset H$ и $e_\ell H=H\cap\Z_\ell=:H_\ell$ (если $h\in H\cap\Z_\ell$, то
$h=e_\ell h$)~--- замкнутая подгруппа $\Z_\ell$. Ясно, что $H\subset\prod H_\ell$. Обратно, элемент $(x_\ell)\in\prod
H_\ell$~--- предел конечных сумм $\sum_{\ell\le L}x_\ell\in H$, а $H$ замкнута.

(c) Если $(\Zh:H)=n$, то $\Zh/H$~--- абелева группа порядка $n$, и $nx\in H$ для всех $x$, т.\,е.\ $n\Zh\subset H$. Подгруппа
$n\Zh=\ker(\Zh\to\Z/n)$ открыта, значит, открыта и $H$ (7.1a(e)); при этом $H$~--- прообраз подгруппы $d\Z/n\Z$
циклической группы $\Z/n$ для некоторого $d\mid n$, т.\,е.\ $H=d\Zh$.
\end{proof}

Удобно кодировать замкнутые подгруппы $\Zh$ \emph{сверхнатуральными числами} (числами Штейница)
$N=\prod_\ell\ell^{n_\ell}$, $n_\ell\in\{0,1,2,\dots,\infty\}$: числу $N$ отвечает $N\Zh=\prod_\ell\ell^{n_\ell}\Z_\ell$
(где $\ell^\infty\Z_\ell=0$). Открытые подгруппы~--- в точности $N\in\mathbb N$ (все $n_\ell<\infty$ и почти все равны $0$).

\begin{corollary}[7.34a]
Подполя $\Fpb$ взаимно однозначно соответствуют сверхнатуральным числам: числу $N$ отвечает
$\F_{p^N}=\bigcup_{n\mid N,\ n<\infty}\F_{p^n}$, и $\Gal(\Fpb/\F_{p^N})=N\Zh$ (при отождествлении 7.17). Каждое
алгебраическое расширение $\F_p$~--- расширение Галуа с проциклической (т.\,е.\ топологически циклической) группой $\Zh/N\Zh\cong\prod_\ell\Z_\ell/\ell^{n_\ell}\Z_\ell$.
\end{corollary}

\begin{proof}
Это 7.13 и 7.34(b): подгруппа $N\Zh$ соответствует полю, неподвижному относительно неё, а это объединение тех
$\F_{p^n}$, на которых она действует тривиально, т.\,е.\ $n\mid N$. Все подгруппы абелевой группы нормальны, поэтому все
подрасширения Галуа (7.13(d)).
\end{proof}

Например, $\F_{p^{2^\infty}}=\bigcup_k\F_{p^{2^k}}$ имеет $\Gal(\F_{p^{2^\infty}}/\F_p)\cong\Z_2$, а
$\Gal(\Fpb/\F_{p^{2^\infty}})\cong\prod_{\ell\ne2}\Z_\ell$; подполя $\F_{p^{2^\infty}}$~--- это $\F_{p^{2^k}}$ и оно само.
Абсолютная группа Галуа любого конечного поля $\F_q$ изоморфна $\Zh$ (с Фробениусом $x\mapsto x^q$ в роли $1$).

\subsection*{Круговые расширения}

\begin{example}[7.35: поле $\Q(\mu_\infty)$]
Пусть $\mu_n\subset\CC$~--- корни степени $n$ из единицы и $\Q(\mu_\infty)=\bigcup_n\Q(\mu_n)$. По 5.8 и 5.10
\[
\chi_n\colon\Gal(\Q(\mu_n)/\Q)\to(\Z/n)^\times,\qquad\sigma(\zeta)=\zeta^{\chi_n(\sigma)}\ (\zeta\in\mu_n),
\]
~--- изоморфизм. При $n\mid m$ ограничение $\Gal(\Q(\mu_m)/\Q)\to\Gal(\Q(\mu_n)/\Q)$ соответствует редукции
$(\Z/m)^\times\to(\Z/n)^\times$ (если $\zeta_n=\zeta_m^{m/n}$, то $\sigma(\zeta_n)=\zeta_n^{\chi_m(\sigma)}$).
Применяя 7.26a к башне $\Q(\mu_{k!})$, получаем изоморфизм топологических групп (\emph{круговой характер})
\[
\chi\colon\Gal(\Q(\mu_\infty)/\Q)\xrightarrow{\ \sim\ }\varlim(\Z/n)^\times=\Zh^\times,
\]
где $\Zh^\times$~--- группа обратимых элементов кольца $\Zh$ (элемент $(c_n)$ обратим $\iff$ все $c_n$ обратимы).
Наглядно: $\sigma(\zeta)=\zeta^{\chi(\sigma)\bmod n}$ для каждого корня из единицы $\zeta$ порядка $n$.

\begin{itemize}[leftmargin=*]
\item \emph{Непрерывность на всей $\Gal(\Qal/\Q)$.} Базу окрестностей единицы в $\Zh^\times$ образуют подгруппы
$U_n=\ker(\Zh^\times\to(\Z/n)^\times)$, а прообраз $U_n$ при
$\Gal(\Qal/\Q)\to\Gal(\Q(\mu_\infty)/\Q)\xrightarrow{\chi}\Zh^\times$~--- это $\Gal(\Qal/\Q(\mu_n))$, открытая подгруппа.
\item \emph{Разложение по простым.} $\Zh^\times\cong\prod_\ell\Z_\ell^\times$, и соответственно
$\Gal(\Q(\mu_\infty)/\Q)\cong\prod_\ell\Gal(\Q(\mu_{\ell^\infty})/\Q)$, где $\Gal(\Q(\mu_{\ell^\infty})/\Q)\cong\Z_\ell^\times$
(та же башня для степеней $\ell$). Строение $\Z_\ell^\times$ дано в 7.33b.
\item \emph{Пример $\ell=2$.} $\Gal(\Q(\mu_{2^\infty})/\Q)\cong\Z_2^\times=\{\pm1\}\times(1+4\Z_2)$. Элемент $-1$ действует
как комплексное сопряжение ($\zeta\mapsto\zeta^{-1}$), поэтому неподвижное поле $\{\pm1\}$~--- вещественное поле
$\bigcup_k\Q(\cos\frac{2\pi}{2^k})$ с группой Галуа $1+4\Z_2\cong\Z_2$. Открытых подгрупп индекса 2 в
$\Z/2\times\Z_2$ ровно три (ядра трёх ненулевых непрерывных гомоморфизмов в $\Z/2$), и им соответствуют три квадратичных
подполя: $\Q(i)$, $\Q(\sqrt2)$, $\Q(\sqrt{-2})$.
\end{itemize}
\end{example}

\subsection*{Композиты}

\begin{example}[7.36]
\emph{(a) Поле $\Q(\mu_\infty,\sqrt[3]2)$.} Пусть $E=\Q(\mu_\infty)$ и $L=\Q(\sqrt[3]2,\omega)$, $\omega=e^{2\pi i/3}$;
$\Gal(L/\Q)\cong S_3$ (перестановки корней $\sqrt[3]2,\omega\sqrt[3]2,\omega^2\sqrt[3]2$). Подполе $E\cap L$ абелево над
$\Q$ (подполе абелева расширения) и содержит $\omega$; подполя $L$, абелевы над $\Q$,~--- неподвижные поля нормальных
подгрупп с абелевым фактором, т.\,е.\ $S_3$ и $A_3$. Значит, $E\cap L=\Q(\omega)=\Q(\sqrt{-3})$. На $\Q(\omega)$ элемент
$a\in\Zh^\times$ действует как $\omega\mapsto\omega^{a\bmod3}$, а $\pi\in S_3$~--- нетривиально тогда и только тогда, когда
$\pi$ нечётна ($A_3=\Gal(L/\Q(\omega))$). По 7.16b
\[
\Gal(\Q(\mu_\infty,\sqrt[3]2)/\Q)\cong\{(a,\pi)\in\Zh^\times\times S_3\mid a\equiv\sgn\pi\pmod3\}.
\]
(Проверка: комплексное сопряжение меняет местами $\omega\sqrt[3]2$ и $\omega^2\sqrt[3]2$~--- нечётная перестановка~--- и
действует на корнях из единицы как $a=-1$.) По 7.15 $\Gal(\Q(\mu_\infty,\sqrt[3]2)/\Q(\mu_\infty))\cong A_3\cong\Z/3$.

\begin{center}
\begin{tikzpicture}[scale=0.9,every node/.style={font=\small}]
\node (T) at (0,2.6) {$\Q(\mu_\infty,\sqrt[3]2)$};
\node (E) at (-2.2,1.3) {$\Q(\mu_\infty)$};
\node (L) at (2.2,1.3) {$\Q(\sqrt[3]2,\omega)$};
\node (I) at (0,0) {$\Q(\omega)$};
\node (Q) at (0,-1.3) {$\Q$};
\draw (L) -- (T);
\draw[cnew,thick] (I) -- node[below right,font=\footnotesize]{$A_3$} (L);
\draw[cnew,thick] (E) -- node[above left,font=\footnotesize]{$A_3$} (T);
\draw (I) -- node[below left,font=\footnotesize]{$\{a\equiv1\ (\mathrm{mod}\ 3)\}$} (E);
\draw (Q) -- node[right,font=\footnotesize]{$\Z/2$} (I);
\end{tikzpicture}
\end{center}
\figcap{Красные стороны имеют изоморфные группы Галуа (7.15).}

\emph{(b) Башня $\Q(\mu_{2^\infty},2^{1/2^\infty})$.} Пусть $\zeta_{2^n}=e^{2\pi i/2^n}$ (так что
$\zeta_{2^{n+1}}^2=\zeta_{2^n}$), $2^{1/2^n}$~--- положительный вещественный корень, $K_n=\Q(\zeta_{2^n},2^{1/2^n})$ и
$K_\infty=\bigcup_nK_n$. Элемент $\sigma\in\Gal(K_\infty/\Q)$ задаётся парой $(b,a)$:
\[
\sigma(2^{1/2^n})=\zeta_{2^n}^{\,b}\,2^{1/2^n},\qquad\sigma(\zeta_{2^n})=\zeta_{2^n}^{\,a},
\qquad b\in\Z/2^n,\ a\in(\Z/2^n)^\times,
\]
согласованно по $n$, т.\,е.\ $b\in\Z_2$, $a\in\Z_2^\times$. Композиция~--- как у аффинных преобразований $x\mapsto ax+b$:
$(b,a)(b',a')=(b+ab',aa')$. Значит, $\Gal(K_\infty/\Q)$ вкладывается в $\Z_2\rtimes\Z_2^\times$.

\emph{Не все пары реализуются.} При $n\ge3$ имеем $\sqrt2=(2^{1/2^n})^{2^{n-1}}$ и $\sqrt2=\zeta_8+\zeta_8^{-1}$. Первое
даёт $\sigma(\sqrt2)=(-1)^b\sqrt2$, второе~--- $\sigma(\sqrt2)=\varepsilon(a)\sqrt2$, где $\varepsilon(a)=1$ при
$a\equiv\pm1$ и $-1$ при $a\equiv\pm3\pmod8$. Значит, $(-1)^b=\varepsilon(a)$.

\emph{Других ограничений нет.} Пусть $M=\Q(\zeta_{2^n})$, $n\ge3$. По теории Куммера (5.30) $[K_n:M]$~--- порядок класса
$2$ в $M^\times/M^{\times2^n}$. Так как $2^{2^{n-1}}=(\sqrt2)^{2^n}$, этот порядок делит $2^{n-1}$. Если бы он делил
$2^{n-2}$, то $(\sqrt2)^{2^{n-1}}=z^{2^n}$ с $z\in M$, откуда $\sqrt2=\xi z^2$, $\xi\in\mu_{2^{n-1}}$; но $\xi$~--- квадрат
в $M$ (так как $\mu_{2^n}\subset M$), и тогда $\sqrt2=w^2$ с $w\in M$. Поле $\Q(w)\subset M$ абелево над $\Q$, значит,
нормально; $w$~--- корень неприводимого (Эйзенштейн) многочлена $X^4-2$, так что $\Q(w)$ степени 4 содержало бы все его
корни $\pm w,\pm iw$, т.\,е.\ поле разложения $X^4-2$ степени 8~--- противоречие. Итак, $[K_n:\Q]=2^{n-1}\cdot2^{n-1}=2^{2n-2}$,
и это ровно порядок подгруппы пар с $(-1)^b=\varepsilon(a)$ в $\Z/2^n\rtimes(\Z/2^n)^\times$. По 7.26a
\[
\Gal(K_\infty/\Q)\cong\{(b,a)\in\Z_2\rtimes\Z_2^\times\mid(-1)^b=\varepsilon(a)\},
\]
замкнутая подгруппа индекса 2. При $n=3$ получается группа поля разложения $X^8-2$ порядка 16. Подгруппа
$\Gal(K_\infty/\Q(\mu_{2^\infty}))$ состоит из пар $(b,1)$ с чётным $b$, т.\,е.\ изоморфна $2\Z_2\cong\Z_2$: <<один
квадратный корень из 2 уже содержится в круговом поле>>. (Ср.\ с 7.38, где ограничений нет.)
\end{example}

\subsection*{Бесконечная теория Куммера и двойственность}

\begin{example}[7.37]
Пусть $F$ содержит примитивный корень степени $n$ из единицы (так что $\operatorname{char}F\nmid n$) и
$F_n=F[(F^\times)^{1/n}]\subset\Fsep$~--- поле, порождённое корнями степени $n$ из всех элементов $F$. Это объединение всех
конечных абелевых расширений $F$ экспоненты, делящей $n$ (5.30), так что $\Gal(F_n/F)$~--- абелева группа экспоненты,
делящей $n$.

\emph{Утверждение.} Спаривание $\langle a,\sigma\rangle=\sigma(a^{1/n})/a^{1/n}$ задаёт изоморфизм топологических групп
\[
\Gal(F_n/F)\xrightarrow{\ \sim\ }\Hom(F^\times/F^{\times n},\mu_n),\qquad\sigma\mapsto\langle\cdot,\sigma\rangle,
\]
где справа~--- все гомоморфизмы с топологией поточечной сходимости (подпространство произведения
$\prod_{a\in F^\times/F^{\times n}}\mu_n$).

\emph{Доказательство.} Для каждой конечной подгруппы $B\subset F^\times/F^{\times n}$ поле $E_B=F[B^{1/n}]$ конечно, и по
5.30--5.32 $\Gal(E_B/F)\cong\Hom(B,\mu_n)$ согласованно с включениями $B\subset B'$ (ограничение автоморфизмов
соответствует ограничению характеров). Поля $E_B$ конфинальны среди конечных подрасширений $F_n$. По 7.26 (в варианте
7.26a)
\[
\Gal(F_n/F)\cong\varlim_B\Hom(B,\mu_n)=\Hom\Bigl(\bigcup_BB,\ \mu_n\Bigr)=\Hom(F^\times/F^{\times n},\mu_n),
\]
и топологии совпадают: в обоих случаях окрестности задаются значениями на конечном числе элементов $a$.

\emph{Почему двойственны прямая сумма и произведение.} Если $F^\times/F^{\times n}\cong\bigoplus_I\Z/n$ (как для $F=\Q$,
$n=2$: базис~--- классы $-1$ и простых), то гомоморфизм прямой суммы задаётся \emph{произвольным} набором значений на
базисе, и $\Gal(F_n/F)\cong\prod_I\mu_n$. Обратно, \emph{непрерывные} гомоморфизмы $\chi\colon\prod_I\Z/n\to\mu_n$ зависят
лишь от конечного числа координат (открытое ядро содержит окрестность $\{x\mid x_i=0,\ i\in J\}$, $J$ конечно), так что
$\Homc(\Gal(F_n/F),\mu_n)\cong\bigoplus_I\Z/n\cong F^\times/F^{\times n}$. Это частный случай \emph{двойственности
Понтрягина}: дискретная абелева группа и компактная абелева группа (здесь~--- экспоненты, делящей $n$) взаимно
двойственны, если характеры компактной группы брать \emph{непрерывными}. Разрывные характеры (они существуют по лемме
Цорна, 7.29--7.30) в эту двойственность не входят~--- так же, как неоткрытые подгруппы не входят в соответствие Галуа.

\emph{Частные случаи.}
\begin{itemize}[leftmargin=*]
\item $F=\R$, $n=2$: $\R^\times/\R^{\times2}=\{\pm1\}$, $\Gal(\CC/\R)=\Z/2$, $\langle-1,\text{сопряжение}\rangle=-1$.
\item $F=\Q$, $n=2$: $F_2$~--- поле $E$ из 7.29, $\Gal(F_2/\Q)\cong\prod\Z/2$. Открытые подгруппы индекса 2
$\leftrightarrow$ ненулевые классы $d\in\Q^\times/\Q^{\times2}$ $\leftrightarrow$ квадратичные поля $\Q(\sqrt d)$.
\item $F=\F_q$, $n\mid q-1$: $\F_q^\times/\F_q^{\times n}\cong\Z/n$, $F_n=\F_{q^n}$, и
$\langle a,\Frob\rangle=(a^{1/n})^{q-1}=a^{(q-1)/n}$. Невырожденность спаривания~--- это \emph{критерий Эйлера}:
$a\in\F_q^{\times n}\iff a^{(q-1)/n}=1$; при $n=2$, $q=p$ это символ Лежандра $\bigl(\frac ap\bigr)$.
\item $F=\Qp$, $p$ нечётно, $n=2$: $\Qp^\times/\Qp^{\times2}=\{1,u,p,up\}\cong(\Z/2)^2$ ($u$~--- единица, не являющаяся
квадратом по модулю $p$; используется лемма Гензеля), так что у $\Qp$ ровно три квадратичных расширения.
\item $F=\CC(t)$: $\CC[t]$ факториально, его простые~--- $t-c$, а константы суть $n$-е степени; поэтому
$\CC(t)^\times/\CC(t)^{\times n}\cong\bigoplus_{c\in\CC}\Z/n$ и $\Gal(F_n/F)\cong\prod_{c\in\CC}\Z/n$.
\end{itemize}
\end{example}

\subsection*{Расширение констант}

\begin{lemma}[7.38a]
Поле $F$ алгебраически замкнуто в поле рациональных функций $F(t)$: каждый элемент $F(t)$, алгебраический над $F$, лежит в
$F$.
\end{lemma}

\begin{proof}
Пусть $r=f(t)/g(t)\in F(t)\setminus F$. Тогда $t$~--- корень многочлена $P(T)=f(T)-r\,g(T)\in F(r)[T]$, который ненулевой
(иначе $f=rg$ и $r\in F$). Значит, $F(t)$ алгебраично над $F(r)$. Если бы $r$ было алгебраично над $F$, то и $t$ было бы
алгебраично над $F$~--- противоречие.
\end{proof}

\begin{example}[7.38]
\emph{(a)} Применим 7.15 к $E=\Fsep$ и $L=F(t)$: $EL=\Fsep(t)$, а $E\cap L=F$ по лемме 7.38a. Значит, $\Fsep(t)/F(t)$~---
расширение Галуа и
\[
\Gal(\Fsep(t)/F(t))\xrightarrow{\ \sim\ }\Gal(\Fsep/F)
\]
(изоморфизм топологических групп). Поле $\Fsep(t)$~--- нормальное промежуточное поле в сепарабельном замыкании
$F(t)^{\rm sep}$, и по 7.13(d) получаем точную последовательность топологических групп
\[
1\to\Gal(F(t)^{\rm sep}/\Fsep(t))\to\Gal(F(t)^{\rm sep}/F(t))\to\Gal(\Fsep/F)\to1.
\]
Слева~--- <<геометрическая>> группа Галуа (константы сепарабельно замкнуты), справа~--- группа Галуа констант. Для
$F=\F_p$ справа стоит $\Zh$, и Фробениус на константах играет центральную роль в арифметике кривых над конечными полями.

\begin{center}
\begin{tikzpicture}[scale=0.9,every node/.style={font=\small}]
\node (S) at (0,3.0) {$F(t)^{\rm sep}$};
\node (Ft) at (0,1.6) {$\Fsep(t)$};
\node (A) at (-2,0.3) {$F(t)$};
\node (B) at (2,0.3) {$\Fsep$};
\node (F) at (0,-1.0) {$F$};
\draw[cold,thick] (Ft) -- node[right,font=\footnotesize,align=left]{геометрическая\\ группа Галуа} (S);
\draw[cnew,thick] (A) -- node[above left,font=\footnotesize]{$\cong\Gal(\Fsep/F)$} (Ft);
\draw (B) -- (Ft); \draw (F) -- (A);
\draw[cnew,thick] (F) -- node[below right,font=\footnotesize]{$\Gal(\Fsep/F)$} (B);
\end{tikzpicture}
\end{center}

\emph{(b) Поле $\Om=\Q(t)(\mu_\infty,t^{1/\infty})$}, порождённое всеми корнями из единицы и согласованными корнями
$t^{1/n}$ ($(t^{1/nm})^m=t^{1/n}$); корни из единицы тоже выбраны согласованно: $\zeta_n=e^{2\pi i/n}$, так что
$\zeta_{nm}^m=\zeta_n$. Как в 7.36(b), элемент группы Галуа задаётся парой $(b,a)$, $\sigma(t^{1/n})=
\zeta_n^bt^{1/n}$, $\sigma(\zeta_n)=\zeta_n^a$. Здесь реализуются \emph{все} пары: по Куммеру
$[\Q(\mu_n)(t)(t^{1/n}):\Q(\mu_n)(t)]$~--- порядок класса $t$ в $\Q(\mu_n)(t)^\times/(\cdot)^{\times n}$, а $t^d$ при
$0<d<n$ не является $n$-й степенью (кратность нуля в $t=0$ у $t^d$ равна $d$, а у $n$-й степени делится на $n$); кроме того,
$\Gal(\Q(\mu_n)(t)/\Q(t))\cong(\Z/n)^\times$ по 7.15 и лемме 7.38a. Степень уровня $n$ равна $n\cdot|(\Z/n)^\times|$, и
\[
\Gal(\Om/\Q(t))\cong\Zh\rtimes\Zh^\times=\{x\mapsto ax+b\mid a\in\Zh^\times,\ b\in\Zh\},
\]
где $\Zh^\times$ действует на $\Zh$ умножением, т.\,е.\ через круговой характер. В отличие от 7.36(b), трансцендентное $t$
не связано никаким соотношением с корнями из единицы, и группа Галуа~--- всё полупрямое произведение.
\end{example}

\begin{aside}[7.39: взгляд вперёд~--- теория полей классов]
\emph{Локально.} Расширения $\Qp$ устроены как расширения $\F_p$ плюс <<ветвление>>. Максимальное неразветвлённое
расширение $\Qp^{\rm ur}=\bigcup_{p\nmid n}\Qp(\mu_n)$ имеет группу $\Gal(\Qp^{\rm ur}/\Qp)\cong\Gal(\Fpb/\F_p)\cong\Zh$ (автоморфизм
определяется действием на поле вычетов; ключ~--- лемма Гензеля), вполне разветвлённые круговые $\Qp(\mu_{p^\infty})$~---
группу $\Zp^\times$, а максимальное абелево расширение $\Qp^{\rm ab}$~--- их композит с группой $\Zh\times\Zp^\times$.
Отображение взаимности $\Qp^\times=p^{\Z}\times\Zp^\times\to\Zh\times\Zp^\times$ непрерывно, инъективно и имеет плотный
образ~--- та же картина <<$\Z$ плотно в $\Zh$>>, что в 7.17: группа Галуа~--- \emph{проконечное пополнение}
мультипликативной группы поля.

\emph{Глобально.} $\Gal(\Qab/\Q)\cong\Zh^\times=\prod_p\Zp^\times$; множитель $\Zp^\times$~--- <<группа инерции в $p$>>. На
$\Q(\mu_n)$ при $\ell\nmid n$ элемент Фробениуса в простом $\ell$~--- это $\zeta\mapsto\zeta^\ell$, т.\,е.\ класс
$\ell\in(\Z/n)^\times$ (на всём $\Q(\mu_\infty)$ он определён лишь с точностью до группы инерции $\Z_\ell^\times$). Теорема
Дирихле о простых в арифметических прогрессиях равносильна тому, что эти классы при каждом $n$ исчерпывают $(\Z/n)^\times$;
теорема Чеботарёва утверждает даже их равнораспределённость. Для произвольного числового поля $K$ группу
$\Gal(K^{\rm ab}/K)$ описывают через \emph{иделы}~--- ограниченное произведение мультипликативных групп $K_v^\times$ всех
пополнений $K$; $\Zh^\times$~--- тень
этого описания для $K=\Q$.
\end{aside}
